%------------------------------------------------------------------------------%
\documentclass[fleqn,utf8,aspectratio=169,12pt,ignorenonframetext]{beamer}

\usepackage{calculo_varias_variaveis-1}

\title{Limites de Funções de Várias Variáveis}

%------------------------------------------------------------------------------%
\begin{document}

\addtitle

%------------------------------------------------------------------------------%
\section{Limites}

%------------------------------------------------------------------------------%
\begin{frame}{Definição}
  Uma função
  \(f\colon R \subset \R^2 \to \R\), \;
  \pause
  \(f(x, y)\),

  \pause
  se aproxima do \emph{limite} $L$
  a medida que \((x, y)\) se aproxima de \((x_0, y_0)\)

  \pause
  se, para todo \emph{$\varepsilon>0$,}
  \pause
  existe um \emph{\(\delta(\varepsilon)>0\)}
  \pause
  tal que,
  \pause
  para todo \((x, y)\) no domínio de $f$

  \pause
  \[
    \abs{f(x, y)-L} < \varepsilon
    \quad\text{ sempre que }\quad
    0 < \sqrt{ \left(x-x_0\right)^2 + \left(y-y_0\right)^2 } < \delta(\varepsilon)
  \]

  \pause
  Obs: \((x_0, y_0)\) não precisa estar no domínio de $f$
\end{frame}

%------------------------------------------------------------------------------%
\begin{frame}{Notação}
  \[
    \lim_{(x, y)\to(x_0, y_0)} f(x, y) = L
  \]
  \pause
  \vspace{\baselineskip}
  \[
    f(x, y) \to L \quad\text{ quando }\quad (x, y)\to(x_0, y_0)
  \]
\end{frame}

%------------------------------------------------------------------------------%
\section{Propriedades}

%------------------------------------------------------------------------------%
\begin{frame}{Teorema -- Unicidade do Limite}
  Se uma função \(f(x, y)\) possui um limite $L$

  quando \((x, y)\) se aproxima de \((x_0, y_0)\),

  \pause
  então \emph{$L$ é único}
\end{frame}

%------------------------------------------------------------------------------%
\begin{frame}{Teorema -- Propriedades dos limites de várias variáveis}
  As propriedades a seguir são verdadeiras se
  \begin{itemize}[<+->]
    \setlength\itemsep{1.5em}
    \item $L$, $M$ e $k$ forem números reais
    \item \(\lim_{(x, y)\to(x_0, y_0)} f(x, y) = L\)
    \item \(\lim_{(x, y)\to(x_0, y_0)} g(x, y) = M\)
  \end{itemize}
\end{frame}

%------------------------------------------------------------------------------%
\begin{frame}{Regra da Soma e Diferença}
  \[
    \lim_{(x, y)\to(x_0, y_0)} \left[f(x, y) + g(x, y)\right]
    \pause =
    \lim_{(x, y)\to(x_0, y_0)} f(x, y) + \lim_{(x, y)\to(x_0, y_0)} g(x, y)
    \pause =
    L + M
  \]
  \pause
  \vspace{\baselineskip}
  \[
    \lim_{(x, y)\to(x_0, y_0)} \left[f(x, y) - g(x, y)\right]
    \pause =
    \lim_{(x, y)\to(x_0, y_0)} f(x, y) - \lim_{(x, y)\to(x_0, y_0)} g(x, y)
    \pause =
    L - M
  \]
\end{frame}

%------------------------------------------------------------------------------%
\begin{frame}{Regra da Multiplicação por Constante e do Produto}
  \[
    \lim_{(x, y)\to(x_0, y_0)} \left[k f(x, y)\right]
    \pause =
    k \lim_{(x, y)\to(x_0, y_0)} f(x, y)
    \pause =
    k L
  \]
  \pause
  \vspace{\baselineskip}
  \[
    \lim_{(x, y)\to(x_0, y_0)} \left[f(x, y) g(x, y)\right]
    \pause =
    \left[\lim_{(x, y)\to(x_0, y_0)} f(x, y) \right]
    \left[\lim_{(x, y)\to(x_0, y_0)} g(x, y) \right]
    \pause =
    LM
  \]
\end{frame}

%------------------------------------------------------------------------------%
\begin{frame}{Regra do Quociente (Divisão)}
  \[
    \lim_{(x, y)\to(x_0, y_0)} \left[ \,\frac{\,f(x, y)\,}{g(x, y)} \, \right]
    \pause =
    \dfrac{\displaystyle\lim_{(x, y)\to(x_0, y_0)} f(x, y)\,}
          {\displaystyle\lim_{(x, y)\to(x_0, y_0)} g(x, y)\,}
    \pause =
    \frac{L}{M}
  \]

  \pause
  \emph{Desde que \(M\neq 0\)}
\end{frame}

%------------------------------------------------------------------------------%
\begin{frame}{Regra da Potência e da Raiz}
  Para $n$ inteiro positivo
  \pause
  \[
    \lim_{(x, y)\to(x_0, y_0)} \left[f(x, y)\right]^n
    \pause =
    \left[\lim_{(x, y)\to(x_0, y_0)} f(x, y)\right]^n
    \pause =
    L^n
  \]
  \pause
  \vspace{\baselineskip}
  \[
    \lim_{(x, y)\to(x_0, y_0)} \sqrt[n]{f(x, y)}
    \pause =
    \sqrt[n]{\lim_{(x, y)\to(x_0, y_0)} f(x, y)}
    \pause =
    \sqrt[n]{L}
  \]

  \pause
  Se $n$ for par $L$ precisa ser positivo
\end{frame}

%------------------------------------------------------------------------------%
\section{Exemplos}

% 01
%------------------------------------------------------------------------------%
\begin{question}
\begin{columns}[t]
\column{0.5\linewidth}
  Calcule os limites
  \vspace{0.5\baselineskip}
  \begin{enumerate}
    \setlength\itemsep{1.5em}
    \item \(\lim_{(x, y)\to(a, b)} k\)
    \item \(\lim_{(x, y)\to(a, b)} x\)
    \item \(\lim_{(x, y)\to(a, b)} y\)
    \item \(\lim_{(x, y)\to(a, b)} xy\)
  \end{enumerate}
\column{0.5\linewidth}
  \pause
  Respostas
  \pause
  \vspace{0.5\baselineskip}
  \begin{enumerate}[<+->]
    \setlength\itemsep{1.5em}
    \item \(\lim_{(x, y)\to(a, b)} k  = k  \)
    \item \(\lim_{(x, y)\to(a, b)} x  = a  \)
    \item \(\lim_{(x, y)\to(a, b)} y  = b  \)
    \item \(\lim_{(x, y)\to(a, b)} xy = ab \)
  \end{enumerate}
\end{columns}
\end{question}

% 02
%------------------------------------------------------------------------------%
\begin{question}
  Calcule o limite
  \(
    \lim_{(x, y)\to(0, 1)} \frac{x-xy+3}{x^2y +5xy - y^3}
  \),
  se ele existir

  \pause
  \vspace{2\baselineskip}
  Assumindo que todos os limites necessários para aplicar as propriedades existam

  (Isso pode não ser verdade)
\end{question}

%------------------------------------------------------------------------------%
\begin{resolution}{Solução}
\begin{align*}
  \onslide<+->{\lim_{(x, y)\to(0, 1)} \frac{x-xy+3}{x^2y +5xy - y^3}}
  \onslide<+->{& = \frac{\displaystyle \lim_{(x, y)\to(0, 1)}(x-xy+3)}
           {\displaystyle \lim_{(x, y)\to(0, 1)}(x^2y +5xy - y^3)} \\}
  \onslide<+->{& = \frac{\displaystyle \lim_{(x, y)\to(0, 1)} x -
                          \lim_{(x, y)\to(0, 1)} xy +
                          \lim_{(x, y)\to(0, 1)} 3}
           {\displaystyle \lim_{(x, y)\to(0, 1)} x^2y +
                          \lim_{(x, y)\to(0, 1)} 5xy -
                          \lim_{(x, y)\to(0, 1)} y^3} \\}
  \onslide<+->{& = \frac{0 - 0\times 1 + 3}{0^2\times 1 + 5\times 0 \times 1 - 1^3} \\}
  \onslide<+->{& = \frac{3}{-1}}
  \onslide<+->{= -3}
\end{align*}
\end{resolution}

% 03
%------------------------------------------------------------------------------%
\begin{question}
  Calcule o limite
  \(
    \lim_{(x, y)\to(0, 0)} \frac{x^2-xy}{\,\sqrt{x}-\sqrt{y}\,}
  \)
  se ele existir

  \pause
  \vspace{\baselineskip}
  Assumindo que todos os limites necessários para aplicar as propriedades existam

  (Isso pode não ser verdade)
\end{question}

%------------------------------------------------------------------------------%
\begin{resolution}{Solução}
\begin{align*}
  \onslide<+->{\lim_{(x, y)\to(0, 0)} \frac{x^2-xy}{\,\sqrt{x}-\sqrt{y}\,}
  & = \frac{\displaystyle \lim_{(x, y)\to(0, 0)} \left(x^2-xy\right)}
           {\displaystyle \lim_{(x, y)\to(0, 0)} \left(\sqrt{x}-\sqrt{y}\right)} \\[2mm]}
  \onslide<+->{& = \frac{\displaystyle \lim_{(x, y)\to(0, 0)} x^2 - \lim_{(x, y)\to(0, 0)}xy}
           {\displaystyle \lim_{(x, y)\to(0, 0)} \sqrt{x} - \lim_{(x, y)\to(0, 0)}\sqrt{y}} \\[2mm]}
  \onslide<+->{& =   \frac{0-0}{0-0} }
  \onslide<+->{\;=\; \frac{0}{0} }
  \onslide<+->{\;=\; 1}
\end{align*}
\onslide<+->{\textcolor{red}{\textbf{Não existe divisão por zero !!!}}}
\end{resolution}

%------------------------------------------------------------------------------%
\begin{resolution}{Solução}
\begin{align*}
  \,\lim_{(x, y)\to(0, 0)} \frac{x^2-xy}{\,\sqrt{x}-\sqrt{y}\,}
  & = \frac{\displaystyle \lim_{(x, y)\to(0, 0)} \left(x^2-xy\right)}
    {\displaystyle \lim_{(x, y)\to(0, 0)} \left(\sqrt{x}-\sqrt{y}\right)} \\[2mm]
  & = \frac{\displaystyle \lim_{(x, y)\to(0, 0)} x^2 - \lim_{(x, y)\to(0, 0)}xy}
    {\displaystyle \lim_{(x, y)\to(0, 0)} \sqrt{x} - \lim_{(x, y)\to(0, 0)}\sqrt{y}} \\[2mm]
  & =  \frac{0-0}{0-0}
  \;=\; \frac{0}{0}
  \qquad\nexists
\end{align*}
\emph{Então o limite não existe?}
\end{resolution}

%------------------------------------------------------------------------------%
\begin{resolution}{Solução}
  \begin{align*}
    \,\lim_{(x, y)\to(0, 0)} \frac{x^2-xy}{\,\sqrt{x}-\sqrt{y}\,}
    & \color{red}
    = \frac{\displaystyle \lim_{(x, y)\to(0, 0)} \left(x^2-xy\right)}
    {\displaystyle \lim_{(x, y)\to(0, 0)} \left(\sqrt{x}-\sqrt{y}\right)} \\[2mm]
    & \color{red} = \frac{\displaystyle \lim_{(x, y)\to(0, 0)} x^2 - \lim_{(x, y)\to(0, 0)}xy}
    {\displaystyle \lim_{(x, y)\to(0, 0)} \sqrt{x} - \lim_{(x, y)\to(0, 0)}\sqrt{y}} \\[2mm]
    & \color{red} = \frac{0-0}{0-0}
    \;=\; \frac{0}{0}
  \end{align*}
  Não podemos aplicar a propriedade do quociente \\  \pause
  \emph{Precisamos remover a indeterminação}
\end{resolution}

%------------------------------------------------------------------------------%
\begin{resolution}{Solução}
  \onslide<+->{Ao calcular o limite consideramos apenas os pontos dentro do domínio da função}

  \onslide<+->{Para \(x\neq y\)}

  \begin{align*}
    \onslide<+->{\frac{x^2-xy}{\,\sqrt{x}-\sqrt{y}\,}}
    \onslide<+->{& = \frac{\left(x^2-xy\right) \emph{\left(\sqrt{x}+\sqrt{y}\right)}}
                          {\left(\sqrt{x}-\sqrt{y}\right)\emph{\left(\sqrt{x}+\sqrt{y}\right)}} \\[1mm]}
    \onslide<+->{& = \frac{x\left(\emph{x-y}\right)\left(\sqrt{x}+\sqrt{y}\right)}{\emph{x-y}} \\[1mm]}
    \onslide<+->{& = x\left(\sqrt{x}+\sqrt{y}\right)}
  \end{align*}
\end{resolution}

%------------------------------------------------------------------------------%
\begin{resolution}{Solução}
\onslide<+->{Calculando o limite}

\begin{align*}
  \onslide<+->{\lim_{(x, y)\to(0, 0)} \frac{x^2-xy}{\,\sqrt{x}-\sqrt{y}\,}}
  \onslide<+->{& = \lim_{(x, y)\to(0, 0)} x\left(\sqrt{x}+\sqrt{y}\right) \\[1mm]}
  \onslide<+->{& = \left(\lim_{(x, y)\to(0, 0)} x\right)
    \left[
    \left(\lim_{(x, y)\to(0, 0)}\sqrt{x}\right) +
    \left(\lim_{(x, y)\to(0, 0)}\sqrt{y}\right)\right] \\[2mm]}
  \onslide<+->{& = 0\left(\sqrt{0}+\sqrt{0}\right) \\[2mm]}
  \onslide<+->{& = 0}
\end{align*}
\end{resolution}

% 04
%------------------------------------------------------------------------------%
\begin{question}
  Analise o limite de \(f(x, y) = \frac{y}{x}\)
  quando \((x, y)\)  se aproxima de \((0, 0)\)

  \pause
  \vspace{\baselineskip}
  A reta \(x=0\) não faz parte do domínio de $f$

  \pause
  O ponto \((0, 0)\) não está no domínio de $f$

  \pause
  Não podemos usar as propriedades diretamente

  \pause
  Não existe manipulação que remova a indeterminação
\end{question}

%------------------------------------------------------------------------------%
\begin{resolution}{Solução}
  Analisando $f$ sobre a reta \(y=x\)
  \pause
  \[
    f(x, y) \evalat{y=x}{}
    \pause
    = f(x, x)
    \pause
    = \frac{x}{x}
    \pause
    = 1
    \pause
    \qquad\text{ pois } x=0 \text{ não está no domínio de } f
  \]

  \pause
  Analisando $f$ sobre a reta \(y=0\)
  \pause
  \[
    f(x, y) \evalat{y=0}{}
    \pause
    = f(x, 0)
    \pause
    = \frac{0}{x}
    \pause
    = 0
    \pause
    \qquad\text{ pois } x=0 \text{ não está no domínio de } f
  \]

  \pause
  O que vai acontecer com outas retas da forma \(y=ax\)?
\end{resolution}

%------------------------------------------------------------------------------%
\begin{resolution}{Solução}
  Para qualquer $\delta>0$
  \pause
  sempre vai existir um

  ponto da forma \((x, x)\) e um ponto da forma \((x, 0)\)

  \pause
  dentro da bola de raio $\delta$ de centro em \((0, 0)\)

  \pause
  Nenhum valor $L$ atende as condições da definição de limite

  \pause
  \vspace{\baselineskip}
  \emph{O limite não existe}
\end{resolution}

%------------------------------------------------------------------------------%
\section{Inexistência do Limite}

%------------------------------------------------------------------------------%
\begin{frame}{Teste dos Dois Caminhos}
  Se \(f(x, y)\) tem limites diferentes ao longo de dois caminhos
  diferentes,

  no domínio de $f$,
  \pause
  quando $(x, y)$ se aproxima de $(x_0, y_0)$,

  \pause
  então
  \[
    \lim_{(x, y)\to(x_0, y_0)} f(x, y) \quad \text{ \emph{não existe}}
  \]
\end{frame}

%------------------------------------------------------------------------------%
\section{Exemplos}

% 05
%------------------------------------------------------------------------------%
\begin{question}
  Verifique se a função
  \(
  f(x, y) = \frac{2x^2y}{x^4+y^2}
  \)
  possui limite na origem.

  \pause
  \vspace{\baselineskip}
  O limite na origem leva a uma indeterminação do tipo $\sfrac{0}{0}$
\end{question}

%------------------------------------------------------------------------------%
\begin{resolution}{Solução}
  Fazendo \(y=mx\), para \(m\neq 0\) e \(x\neq 0\),
  \pause
  temos
  \[
    f(x, y)\evalat{y=mx}{}
    \pause
    = \frac{2x^2y}{x^4+y^2} \evalat{y=mx}{}
    \pause
    = \frac{2mx^3}{x^4+m^2x^2}
    \pause
    = \frac{2mx\,x^2}{\left(x^2+m^2\right)x^2}
    \pause
    = \frac{2mx}{x^2+m^2}
  \]

  \pause
  Calculando o limite \(x\to 0\), que corresponde a \((x, y)\to (0, 0)\)
  \[
    \pause
    \lim_{x\to 0} \frac{2mx}{x^2+m^2}
    \pause
    = \frac{0}{m^2}
    \pause
    = 0
    \qquad\qquad
    \forall m
  \]

  \pause
  Portanto,
  \(
    {
      \color<12->{red}
      \lim_{(x, y)\to(0, 0)} \frac{2x^2y}{x^4+y^2} = 0
    }
    \pause
    \qquad
    \qquad
    \text{\textcolor{red}{\textbf{Falso!}}}
  \)
\end{resolution}

%------------------------------------------------------------------------------%
\begin{resolution}{Solução}
  Fazendo \(y=kx^2\), para \(k\neq 0\) e \(x\neq 0\),
  \pause
  temos
  \[
  f(x, y)\evalat{y=kx^2}{}
  \pause
  = \frac{2x^2y}{x^4+y^2} \evalat{y=kx^2}{}
  \pause
  = \frac{2kx^4}{x^4+k^2x^4}
  \pause
  = \frac{2kx^4}{\left(1+k^2\right)x^4}
  \pause
  = \frac{2k}{1+k^2}
  \]

  \pause
  Para cada valor de $k$, a função $f$ se aproxima da origem com um valor diferente

  \pause
  Portanto, o limite não existe
\end{resolution}

%------------------------------------------------------------------------------%
\section{Lista Mínima}

%------------------------------------------------------------------------------%
\begin{frame}{Lista Mínima}
  Cálculo Vol. 2 do Thomas 12\textsuperscript{a} ed. -- Seção 14.2
  \begin{enumerate}
    \item Estudar o texto da seção
    \item Resolver os exercícios: 4-6, 13-15, 18, 26, 55
  \end{enumerate}

  \vfill
  Atenção: A prova é baseada no livro, não nas apresentações
\end{frame}

%------------------------------------------------------------------------------%
\end{document}
%------------------------------------------------------------------------------%
